\section{Arithmétique des lacunes entre les échelles \texorpdfstring{\(3^6\)}{3 élevé à 6} et \texorpdfstring{\(10^3\)}{10 élevé à 3}}
\label{sec:vacancias-corriente-electron-rev11}
\HMTClaim{MD-R-VACANCY-POLARIZATION-CURRENT}
\index{lacune!calendrier arithmétique}
\index{polarisation!courant discret}

Les échelles finies
\[
 729=3^6,
 \qquad
 1000=10^3
\]
déterminent un décalage logarithmique exact. Nous définissons
\begin{equation}
 \rho_{\rm vac}=\log_{1000}729=\log_{10}9,
 \qquad
 \delta_{\rm vac}=1-\rho_{\rm vac}=\log_{10}(10/9).
 \label{eq:delta-vacancias-rev11}
\end{equation}
L'irrationalité de \(\delta_{\rm vac}\) s'obtient directement à partir des
valuations premières : une égalité
\(\delta_{\rm vac}=p/q\) impliquerait
\((10/9)^q=10^p\), incompatible avec les valuations en \(2,3,5\).
La lacune est donc une rotation arithmétique apériodique déterminée
entièrement par le rapport entre les deux échelles ; sa définition précède
la loi électronique.

\subsection{Lois exactes d'espacement et de retour}

Soit
\[
 p_n=\left\lfloor\frac{n}{\delta_{\rm vac}}\right\rfloor.
\]

\begin{theorem}[Loi des espacements \texorpdfstring{\(21/22\)}{21/22}]
\label{thm:ley-huecos-vacancias-rev11}
Les incréments consécutifs de la suite \(p_n\) satisfont
\[
 p_{n+1}-p_n\in\{21,22\}.
\]
Les deux valeurs apparaissent dans le mot mécanique.
\end{theorem}

\begin{proof}
Pour tout irrationnel \(x>1\), la différence
\(\lfloor(n+1)x\rfloor-\lfloor nx\rfloor\) appartient à
\(\{\lfloor x\rfloor,\lceil x\rceil\}\). Dans ce cas
\[
 \delta_{\rm vac}^{-1}
 =21.85434532678283256\ldots,
\]
ce qui détermine les deux espacements. L'irrationalité exclut la stabilisation
sur un seul d'entre eux.
\end{proof}

La fréquence de l'espacement court est
\[
 \sigma_{\rm vac}=22-\delta_{\rm vac}^{-1}
 =0.1456546732171674374\ldots,
 \qquad
 \sigma_{\rm vac}^{-1}=6.865553832996660997\ldots.
\]

\begin{corollary}[Loi de retour \texorpdfstring{\(6/7\)}{6/7}]
\label{cor:ley-retorno-vacancias-rev11}
Les occurrences successives de l'espacement \(21\) sont séparées par six ou
sept positions.
\end{corollary}

\begin{proof}
Le mot indicateur de l'espacement court est le codage mécanique de
pente irrationnelle \(\sigma_{\rm vac}\). L'argument de Beatty appliqué à
\(\sigma_{\rm vac}^{-1}\) produit exactement les retours \(6/7\)
\cite{Lothaire2002}.
\end{proof}

Les deux lois forment le cadre orienté
\begin{equation}
 V_{\rm vac}=
 \begin{pmatrix}21&22\\6&7\end{pmatrix},
 \qquad
 \det V_{\rm vac}=21\cdot7-22\cdot6=15.
 \label{eq:marco-vacancias-21-22-rev11}
\end{equation}
Le déterminant non nul prouve l'indépendance des directions
d'espacement et de retour. Le cardinal \(15\) réapparaît ainsi comme invariant du
calendrier des lacunes. Dans l'équation électronique, le couple
\((-22,+15)\) provient du sélecteur central démontré dans
\cref{thm:coeficientes-electron-rev6} ; le repère
\eqref{eq:marco-vacancias-21-22-rev11} apporte une corroboration arithmétique
indépendante, avec sa propre généalogie.

Il convient de distinguer deux mots dérivés du même décalage. La
fonction indicatrice
\[
 h_n=\mathbf 1_{\{p_{n+1}-p_n=21\}}
\]
a pour fréquence \(\sigma_{\rm vac}\) et pour retours \(6/7\). Le mot
\[
 z_n=\lceil(n+1)\delta_{\rm vac}\rceil
      -\lceil n\delta_{\rm vac}\rceil
\]
a pour pente \(\delta_{\rm vac}\) et contient \(45\) ou \(46\) unités
dans chaque bloc de longueur mille. Les lois \(21/22\), \(6/7\) et \(45/46\)
sont des lecteurs distincts d'une même rotation irrationnelle.

\subsection{Polarisation bornée et équation de continuité discrète}

Une phase \(\theta\) définit
\[
 z_n(\theta)=\lceil(n+1)\delta_{\rm vac}+\theta\rceil
 -\lceil n\delta_{\rm vac}+\theta\rceil.
\]
La polarisation intrinsèque accumulée est
\begin{equation}
 \begin{aligned}
 P_N(\theta)
 &:=\sum_{n=0}^{N-1}\bigl(z_n(\theta)-\delta_{\rm vac}\bigr)\\
 &=\lceil N\delta_{\rm vac}+\theta\rceil-
   \lceil\theta\rceil-N\delta_{\rm vac},
 \qquad |P_N(\theta)|<1.
 \end{aligned}
 \label{eq:polarizacion-vacancias-rev11}
\end{equation}
La différence temporelle
\begin{equation}
 J_N(\theta)=P_{N+1}(\theta)-P_N(\theta)
 =z_N(\theta)-\delta_{\rm vac}
 \label{eq:corriente-vacancias-rev11}
\end{equation}
est le courant arithmétique associé. Les
\cref{eq:polarizacion-vacancias-rev11,eq:corriente-vacancias-rev11}
constituent une équation de continuité exacte : la charge accumulée demeure
uniformément bornée et chaque impulsion est sa variation locale. La réalisation
électromagnétique applique à ce couple une unité de charge et une échelle temporelle ;
la structure de polarisation et de courant est déjà fixée avant cette carte.

\subsection{Clôture harmonique de modes nuls conjugués dans l'état électronique}
\label{sec:nido-armonico-modos-nulos-rev11}
\HMTClaim{MD-R-ELECTRON-NULL-MODE-NEST-CANDIDATE}

La section électronique centrale admet une lecture structurelle comme clôture
de deux modes nuls conjugués. Pour \(u,v>0\), les projecteurs hyperboliques
\(P_\pm\) satisfont
\[
 N_{\rm sp}(uP_++vP_-)=uv>0,
\]
bien que chaque composante isolée ait une norme nulle. Dans cette construction, la
clôture résonante de modes nuls électromagnétiques conjugués désigne l'anneau
de transport à deux feuillets dont l'holonomie ferme la section matérielle
centrale ; l'émission et l'absorption correspondent à l'ouverture et à la fermeture d'une
orientation nulle par rapport à cette section.

\begin{statusbox}[title={Réalisation photonique de la clôture harmonique}]
La clôture de deux modes nuls et la positivité de leur norme couplée sont des
identités algébriques. La réalisation électromagnétique est spécifiée par
un opérateur d'entrelacement avec les modes de Maxwell, un hamiltonien de clôture et les
taux de transition correspondants. L'expression publique
\emph{clôture résonante de modes nuls électromagnétiques conjugués}
nomme l'objet structurel qui reçoit cette réalisation.
\end{statusbox}
